\begin{ouverture}[college3]{14}{Loi d’Ohm et puissance électrique}{Partie 3 · L’énergie, ses transferts et ses conversions}{Comment prévoir le fonctionnement d’un dipôle ?}
\objectif{Exploiter U = RI et une caractéristique.}
\objectif{Calculer P = UI en courant continu.}
\objectif{Contrôler unités, limites et puissance admissible.}
\prerequis{Lois des circuits et montage A/V.}
\prerequis{Proportionnalité.}
\end{ouverture}
\pcPage{Le cours}{Loi d’Ohm et puissance électrique}
\phantomsection\label{ess-14}
\essSection{1}{Un conducteur ohmique}
Pour un conducteur ohmique dans des conditions physiques fixées, notamment de température, la tension U est proportionnelle à l’intensité I : $U=RI$. U en V, I en A et R en ohms, $\Omega$. La caractéristique U en fonction de I est une droite passant par l’origine dans le modèle. La pente U/I vaut R. Tous les dipôles ne sont pas ohmiques.\par
\pcFigure{figures/b14-ohm.pdf}{Résistance de 100 Ω : attention à l’unité de l’axe des intensités.}
\essSection{2}{Isoler l’inconnue}
On peut écrire $I=U/R$ ou $R=U/I$ pour I non nul. Convertir les mA en A et les k$\Omega$ en $\Omega$. À tension fixée, une résistance plus grande laisse passer une intensité plus faible. Une lampe peut chauffer et présenter une caractéristique différente : on ne lui applique pas automatiquement une résistance constante.\par
\clearpage
\essSection{3}{Puissance électrique}
En courant continu, la puissance électrique reçue par un dipôle vaut $P=UI$, avec P en W, U en V et I en A. La puissance est une énergie transférée par unité de durée. Une puissance nominale décrit un fonctionnement prévu dans les conditions indiquées. Pour un conducteur ohmique, combiner U = RI et P = UI permet de vérifier les risques d’échauffement.\par
\pcFigure{figures/b14-montage.pdf}{Mesurer U et I avant de calculer P ; ne pas dépasser les caractéristiques du matériel.}
\essSection{4}{Vérifier un choix}
Une résistance possède une puissance maximale admissible donnée par le fabricant. Un calcul cohérent de R ne garantit pas que l’échauffement soit acceptable. Les montages du collège utilisent une alimentation scolaire adaptée ; les installations domestiques demandent d’autres modèles et des protections. P = UI n’est utilisé ici qu’en courant continu ou dans les cas explicitement simplifiés par l’énoncé.\par
\begin{pcEncadre}{method}{Un exemple expliqué}Une résistance de $200\unite\Omega$ reçoit $6\unite V$. $I=6/200=0{,}030\unite A=30\unite{mA}$. Sa puissance vaut $P=6\times0{,}030=0{,}18\unite W$. Elle reste inférieure à une limite fournie de 0,25 W ; le professeur vérifie aussi les autres contraintes du montage.\end{pcEncadre}
\pcPage{Expérimenter et retenir}{Loi d’Ohm et puissance électrique}
\phantomsection\label{activite-14}
\begin{pcEncadre}{method}{Investigation en classe}Relever U et I pour une résistance avec une alimentation scolaire réglable, sans dépasser les limites données. Faire un nuage de points U(I) avec I en A. Comparer les quotients hors origine et l’alignement, puis calculer la puissance au point le plus élevé. Signaler toute dérive due à l’échauffement.\end{pcEncadre}
\pcSousTitre{Ce que je dois savoir expliquer}\begin{questions}
\item Reconnaître et exploiter une proportionnalité U/I.
\item Écrire et convertir avant de calculer.
\item Comparer la puissance calculée à une limite fournie.
\end{questions}
\begin{pcEncadre}{attention}{Une erreur à éviter}Si l’axe est en mA, sa valeur numérique n’est pas en A : diviser directement 6 V par 30 donne une erreur d’un facteur 1 000 sur la résistance.\end{pcEncadre}
\pcRepere{cor-activite-14}{Pistes pour l’investigation}\pcGrille
\clearpage\markright{Je démarre}
\rubrique{Je démarre}
\par\noindent\begin{minipage}{\linewidth}
\begin{exercice}[1]{Les symboles}
\begin{questions}
\item Dans U = RI, que désignent U et R ?
\item Donner leurs unités.
\end{questions}
\end{exercice}
\end{minipage}
\par\noindent\begin{minipage}{\linewidth}
\begin{exercice}[1]{Convertir}
\begin{questions}
\item Convertir $2{,}2\unite{k\Omega}$ en $\Omega$.
\item Convertir 15 mA en A.
\end{questions}
\end{exercice}
\end{minipage}
\par\noindent\begin{minipage}{\linewidth}
\begin{exercice}[1]{Trouver U}
\begin{questions}
\item Calculer U pour $R=100\unite\Omega$ et $I=0{,}020\unite A$.
\item Donner la relation utilisée.
\end{questions}
\end{exercice}
\end{minipage}
\par\noindent\begin{minipage}{\linewidth}
\begin{exercice}[1]{Trouver I}
\begin{questions}
\item Calculer I sous 6 V pour $R=300\unite\Omega$.
\item Exprimer en mA.
\end{questions}
\end{exercice}
\end{minipage}
\par\noindent\begin{minipage}{\linewidth}
\begin{exercice}[1]{Trouver R}
\begin{questions}
\item Un dipôle ohmique reçoit 3 V sous 10 mA. Calculer R.
\item Peut-on trouver R par U/I à l’origine où U = I = 0 ?
\end{questions}
\end{exercice}
\end{minipage}
\clearpage\markright{Je m’entraîne}
\rubrique{Je m’entraîne}
\par\noindent\begin{minipage}{\linewidth}
\begin{exercice}[2]{Une puissance}
\begin{questions}
\item Une lampe en continu reçoit 12 V sous 0,50 A. Calculer P.
\item Quelle énergie reçoit-elle chaque seconde ?
\end{questions}
\end{exercice}
\end{minipage}
\par\noindent\begin{minipage}{\linewidth}
\begin{exercice}[2]{Retrouver l’intensité}
\begin{questions}
\item Un appareil continu de 24 W fonctionne sous 12 V. Trouver I.
\item Un second appareil de 12 W sous 12 V demande quelle intensité ?
\end{questions}
\end{exercice}
\end{minipage}
\par\noindent\begin{minipage}{\linewidth}
\begin{exercice}[2]{Une caractéristique}
\pcFigureEx{figures/rev-ex-ohm-100.pdf}{}
\begin{questions}
\item On relève (I en A ; U en V) : (0,01 ; 1), (0,02 ; 2), (0,03 ; 3). Trouver R.
\item Quel contrôle graphique effectuer ?
\end{questions}
\end{exercice}
\end{minipage}
\par\noindent\begin{minipage}{\linewidth}
\begin{exercice}[2]{Une limite dépassée}
\begin{questions}
\item Une résistance de $100\unite\Omega$ sous 10 V est limitée à 0,25 W. Calculer I puis P.
\item Le choix respecte-t-il la limite ?
\end{questions}
\end{exercice}
\end{minipage}
\par\noindent\begin{minipage}{\linewidth}
\begin{exercice}[2]{À tension constante}
\begin{questions}
\item R double sous la même tension. Que devient I ?
\item Que devient P = UI ?
\end{questions}
\end{exercice}
\end{minipage}
\clearpage\markright{Je raisonne}
\rubrique{Je raisonne}
\par\noindent\begin{minipage}{\linewidth}
\begin{exercice}[3]{À résistance constante}
\begin{questions}
\item U double et R reste constante. Comment évolue I ?
\item Comment évolue P ?
\end{questions}
\end{exercice}
\end{minipage}
\par\noindent\begin{minipage}{\linewidth}
\begin{exercice}[3]{Une lampe réelle}
\begin{questions}
\item Une lampe donne U/I = 20 Ω puis 35 Ω à deux fonctionnements. Est-elle ohmique à R constante sur cet intervalle ?
\item Quelle cause physique peut contribuer à cette variation ?
\end{questions}
\end{exercice}
\end{minipage}
\par\noindent\begin{minipage}{\linewidth}
\begin{exercice}[3]{Un graphique trompeur}
\pcFigureEx{figures/rev-ex-ohm-inverse.pdf}{Exemple avec la tension en abscisse et l’intensité en ordonnée.}
\begin{questions}
\item Un graphe porte U horizontalement et I verticalement. Sa pente vaut-elle R ?
\item Comment retrouver R sur un point hors origine ?
\end{questions}
\end{exercice}
\end{minipage}
\par\noindent\begin{minipage}{\linewidth}
\begin{exercice}[3]{Deux branches}
\begin{questions}
\item Deux résistances de 100 Ω et 200 Ω sont en dérivation sous 6 V. Trouver leurs courants.
\item Calculer la puissance totale reçue.
\end{questions}
\end{exercice}
\end{minipage}
\par\noindent\begin{minipage}{\linewidth}
\begin{exercice}[3]{Choisir un dipôle}
\begin{questions}
\item Un conducteur doit être parcouru par 20 mA sous 5 V. Trouver R et P.
\item Une limite de 0,125 W convient-elle dans le modèle ? Discuter la marge.
\end{questions}
\end{exercice}
\end{minipage}
\pcPage{Défi du chapitre}{Loi d’Ohm et puissance électrique}
\phantomsection\label{defi-ex-14}
\pcDefiTitre{14}{Une résistance pour le laboratoire}
\begin{questions}
\item On dispose de résistances de 100 Ω ou 220 Ω, chacune limitée à 0,25 W. La tension continue maximale est 6 V. Calculer I et P pour chaque choix.
\item Quel choix respecte la limite donnée ?
\item Décrire les mesures qui permettraient de contrôler le fonctionnement sans confondre les deux appareils.
\end{questions}
\pcGrille
